\paragraph{Exact behaviour under every constant extension.}
Keep this particular degree-nineteen cover over $\mathbf F_{64}$
and its Prym $B$.  For $m\geq1$, let $M_m$ be the exponent of
$B(\mathbf F_{64^m})[2^\infty]$, with exponent one for the zero
group.  Then
\[
 M_m=
 \begin{cases}
  1,&7\nmid m,\\
  2,&7\mid m\text{ and }m\text{ is odd},\\
  2^{v_2(m)+2},&14\mid m.
 \end{cases}                                                   \tag{BE.22}
\]
Consequently this same cover satisfies A34's sufficient inequality
\[
                    19>(2M_m-1)^2
                    \quad\Longleftrightarrow\quad14\nmid m.   \tag{BE.23}
\]
For every field in the successful cases it separates the section
of $W$ from every distinct rational point.  Failure of this
inequality when $14\mid m$ does not assert failure of point
separation by this cover or by other covers.

\begin{proof}
Use the already checked character polynomial
$L_\chi(T)=1+cT+64T^2$ and the complete unramified ring
$\mathcal O=\mathbf Z_2[\zeta_{19}]$ of degree eighteen.
By (BE.19), $c$ is a unit.  The reduction of
$X^2+cX+64$ has the two distinct roots $0$ and $\bar c$.
Hensel's lemma for complete local rings
\cite[Lemma~10.153.9, Tag~\texttt{04GM}]{Stacks04GM}
gives a unit root $u\in\mathcal O$; its other root is
$w=64/u$.  The root equation implies
\[
 u\equiv-c\pmod{64\mathcal O},\qquad
 \operatorname{ord}(\bar u)=7,
 \qquad u^7\equiv3\pmod{8\mathcal O}.                          \tag{BE.24}
\]
For the last congruence use $1+c^7\equiv6\pmod8$ from
(BE.19), so $(-c)^7\equiv3\pmod8$.

The deck action is already constant over $\mathbf F_{64}$.
Thus on the same faithful character space, extension of
constants to degree $m$ raises both Frobenius roots to their
$m$th powers.  The value of the resulting character factor is
\[
           L_{\chi,m}(1)=(1-u^m)(1-w^m).
                                                               \tag{BE.25}
\]
Since $v_2(w)=6$, the second factor is a unit.  Formula
(BE.8) applies over each stated extension and gives
\[
 \begin{gathered}
 B(\mathbf F_{64^m})[2^\infty]\simeq
      \mathcal O/2^{e_m}\mathcal O,
 \qquad e_m=v_2(1-u^m),\\
 v_2\bigl(\#B(\mathbf F_{64^m})\bigr)=18e_m.
 \end{gathered}
                                                               \tag{BE.26}
\]
No identification of a crystalline or etale Frobenius
eigenvalue is needed beyond the roots of the established
character polynomial and the finite-module proof of (BE.8).

If $7\nmid m$, (BE.24) gives $e_m=0$.
Otherwise write $m=7j$ and put $a=u^7$.
The congruence $a\equiv3\pmod8$ gives
$v_2(a-1)=1$ and $v_2(a+1)=2$.
For odd $j$, the factor $1+a+\cdots+a^{j-1}$ reduces to
one modulo two, so $v_2(a^j-1)=1$.
For even $j=2^t b$ with $t\geq1$ and $b$ odd, first
$v_2(a^2-1)=3$.  Each subsequent squaring increases this
valuation by exactly one: the previous even power is one
modulo eight, so its sum with one has valuation one.
The final odd power $b$ does not change the valuation,
by the same geometric-sum argument.  Hence
$e_m=t+2=v_2(m)+2$.  This proves (BE.22).

In its first two cases $M_m$ is respectively one or two,
so the right side of A34's inequality is respectively one
or nine.  If $14\mid m$, then $M_m\geq8$ and that right
side is at least $225$.  This proves (BE.23).
\end{proof}
